% ch-a2-06 —— 第 2 篇 Zhu，PDF 117–126，印刷页 104–113
\ysrc{104}
下面给出 $\mathfrak A_{\mathrm{GL}_n,\{1\}}$ 的猜想性描述。考虑分类链 $\{V_1\to V_2\to\cdots\to V_n\}$ 的导出叠 $\mathcal W_n$，其中 $V_i$ 是 $W_F$ 的 $i$ 维表示（即 $V_i\in\Loc_{{}^{c}\mathrm{GL}_i,F}$）。只记住 $V_n$ 就给出自然态射 $\pi:\mathcal W_n\to\Loc_{{}^{c}\mathrm{GL}_n,F}$。于是 Braverman–Finkelberg 猜想的算术类比是
\[
\mathfrak A_{\mathrm{GL}_n,\{1\}}\cong\mathfrak A_{\mathrm{BF}}:=\pi_*\omega_{\mathcal W_n}.
\]
把 $\mathfrak A_{\mathrm{GL}_n,\{1\}}$ 的这两个猜想性描述合起来，就得到下述猜想。

\begin{conj}{4.29}\label{thm:a2-4-29}
作为 $\Loc_{{}^{c}\mathrm{GL}_n,F}$ 上的拟凝聚层，$\mathfrak A_{\mathrm{EH}}$ 与 $\mathfrak A_{\mathrm{BF}}$ 之间存在自然同构。
\end{conj}

\begin{rem}{4.30}\label{thm:a2-4-30}
为讨论 $\mathfrak A_G$ 对 $t$ 的依赖，本注把它写作 $\mathfrak A^t_G$。若 $\theta\in\mathbf{Tor}_{Z_G,\mathrm{iso}F}$ 把 $t_1\in\mathbf{TS}_G$ 送到 $t_2\in\mathbf{TS}_G$，则应存在函子的典范同构
\begin{equation}
\mathfrak A^{t_2}_G(-)\cong\mathfrak A^{t_1}_G(-)\otimes\mathcal L_\theta, \tag{4.16}\label{eq:a2-4-16}
\end{equation}
其中 $\mathcal L_\theta$ 如猜想 3.20 所述。更精确地说，应存在 $\mathbf{Tor}_{Z_G,\mathrm{iso}F}$-等变的正合全忠实函子
\[
\mathfrak A_G:\Rep_{\mathrm{f.g.}}(G(F),\Lambda)\times^{\mathbf{Tor}^0_{Z_G}}\mathbf{TS}_G\to\Coh(\Loc_{{}^{c}G,F,\iota}).
\]
若 $G$ 拟分裂，则 $\mathfrak A_G$ 由下述典范全忠实函子诱导得到：
\[
\Rep_{\mathrm{f.g.}}(G(F),\Lambda)\times^{\mathbf{Tor}^0_{Z_G}}\mathbf W_G\to\Coh(\Loc_{{}^{c}G,F,\iota}).
\]
\end{rem}

下面记下该猜想的一个推论。回忆稳定中心 $Z^{c}_{G,F}$（见 (3.5)），以及 $G$ 的 Hecke 代数 $H_G$（见注 4.3）。令 $Z_{G,F}:=Z(H_G)$ 记 $H_G$ 的中心，即 $G(F)$ 的 Bernstein 中心。

\begin{cor}{4.31}\label{thm:a2-4-31}
假定上述猜想成立，则存在自然映射
\begin{equation}
Z^{c}_{G,F}\to Z_{G,F}, \tag{4.17}\label{eq:a2-4-17}
\end{equation}
它不依赖 $t\in\mathbf{TS}_G$ 的选取。此外，该映射应与抛物诱导相容（这将特别蕴含 (4.14)）。对 $\Loc_{{}^{c}G,F,\iota}$ 的连通分支 $D$，令 $Z^{c}_{G,F,D}$ 与 $Z_{G,F,D}$ 为相应的幂等分量。则 $Z_{G,F,D}$ 在 $Z^{c}_{G,F,D}$ 上有限。若 $G=G^*$，则 \eqref{eq:a2-4-17} 是分裂单射。
\end{cor}

\begin{rem}{4.32}\label{thm:a2-4-32}
当 $G=\mathrm{GL}_n$、$F$ 为 $p$-进域、$\Lambda=\overline\Q$ 时，推论中的映射已由 Scholze [65] 构造。利用 $\mathrm{GL}_n$ 的局部 Langlands 对应，Helm 与 Helm–Moss 对 $\Lambda=\overline\Z_\ell$ 构造了这样的映射 [39–41]。注意 $\mathrm{GL}_n$ 时 (4.17) 是同构。对一般的 $G$，从游历代数（见注 3.15）到 $Z_{G,F}$ 的映射由 Genestier–Lafforgue [36] 构造（在等特征情形以及 $\ell$-进完备化之后）。(4.17) 的一般情形（$\Lambda=\Z_\ell$）出现在 Fargues–Scholze 的工作中 [27]，但那里没有构造 $\mathfrak A_G$。至于 $Z_{G,F}\to Z_{G,F}$ 的有限性（限制到 $\Loc_{{}^{c}G,F,\iota}$ 的每个分支 $D$ 上），最近在 [20] 中得到证明。\tn{底本此处前后两项都写作 $Z_{G,F}$，按推论 4.31 的陈述，前一项应为 $Z^{c}_{G,F}$，疑为底本漏排上标 $c$。}
\end{rem}

\begin{rem}{4.33}\label{thm:a2-4-33}
若 $G=T$ 是环面，则 (4.17) 的存在性应来自猜想 3.18，而后者又应诱导函子
\[
\Rep(T(F),\Lambda)\cong\mathbf{Mod}_{Z^{c}_{T,F}}\subset\Qcoh(\Loc_{{}^{c}T,F,\iota}),
\]
把 $\Rep_{\mathrm{f.g.}}(T(F),\Lambda)$ 送到 $\Coh(\Loc_{{}^{c}T,F})$。这应当就是所要的函子 $\mathfrak A_T$。
\end{rem}

\ysrc{105}
遗憾的是，目前对一般的 $\mathfrak A_{G,K}$ 还得不到显式的猜想性描述。下面是一些预期与注记。

(1) 我们预期：若 $K$ 是某个 parahoric 子群的 pro-unipotent 根，则 $\mathfrak A_{G,K}$ 支撑在 $\Loc^{\mathrm{tame}}_{{}^{c}G,F,\iota}$ 上。特别地，应存在映射
\begin{equation}
Z^{\mathrm{tame}}_{{}^{c}G,F}\to Z(H_{G,K}), \tag{4.18}\label{eq:a2-4-18}
\end{equation}
它是 (4.13) 的推广。

(2) 设 $G$ 拟分裂。我们预期：对一族共尾的开紧子群 $K\subset G(F)$，存在拟光滑导出叠 $\widetilde{\Loc}^{K}_{{}^{c}G,F,\iota}$ 与本征概形态射 $\pi^K:\widetilde{\Loc}^{K}_{{}^{c}G,F,\iota}\to\Loc_{{}^{c}G,F,\iota}$，使得
\[
\mathfrak A_{G,K}\cong\pi^K_*\mathcal O_{\widetilde{\Loc}^{K}_{{}^{c}G,F,\iota}}\cong\pi^K_!\omega_{\widetilde{\Loc}^{K}_{{}^{c}G,F,\iota}}.
\]
注意这将特别蕴含 $\mathfrak A_{G,K}$ 关于 Grothendieck–Serre 对偶自对偶（见注 4.27）。

(3) 利用这样一个事实：$\Loc_{{}^{c}G,F,\iota}$ 的某个连通分支「看起来像」另一个群 $G'$ 的驯顺 Langlands 参数叠（见命题 2.35 的证明），也许可以把 $\mathfrak A_G$ 在该分支上的限制与另一个群的凝聚 Springer 层联系起来。若 $G=\mathrm{GL}_n$，这也许能给出 $\mathfrak A_G$ 的「手工」构造。沿这一思路的进路见 [5]。

(4) 即便对各个 $\Lambda$ 都弄清了 $\{\mathfrak A_{G,K}\}_K$（从而知道函子 $\mathfrak A_G$ 定义合理），理解与具体 $G(F)$-表示相对应的 $\Loc_{{}^{c}G,F,\iota}$ 上 (ind)-凝聚层仍然重要（有时也颇有难度）。举例来说，设 $X$ 是 $F$ 上的 $G$-簇。则 $C_c(X(F))$ 是自然的 $G(F)$-表示，因而应对应 $\Loc_{{}^{c}G,F,\iota}$ 上的 ind-凝聚层 $\mathfrak A_X:=\mathfrak A_G(C_c(X(F)))$。Ben-Zvi–Sakellaridis–Venkatesh 最近在相对 Langlands 纲领中提出的猜想，在当前设定下应当也有类比，（至少当 $\Lambda$ 是特征零的域时）从而对某些 $X$ 纯从伽罗瓦侧给出 $\mathfrak A_X$ 的猜想构造。

\subsection{范畴算术局部 Langlands 对应}\label{sec:a2-4-6}

本小节解释猜想的层 $\mathfrak A_G$ 如何纳入局部 Langlands 猜想的一个假想的范畴形式。更详细的讨论将出现在 [88] 中。设 $\Lambda$ 是 $\Z_\ell$ 上的代数，其中 $\ell\ne p$。为简单起见，本小节把 $\Loc_{{}^{c}G,F}\otimes_{\Z_\ell}k$ 写作 $\Loc_{{}^{c}G}$。固定一个导子为 $\mathcal O_F$ 的非平凡特征 $\psi_0:F\to\Lambda^\times$。

在局部 Langlands 中，许多人有一个共识：最好不要只研究单个 $p$-进群 $G$ 的表示论，而应同时研究一族与 $G$ 密切相关的群的表示论。把这想法精确化的方式有多种，只需适当选取这族群，例如 Vogan 的纯内型、Kottwitz–Kaletha 的扩张纯内型，或 Kaletha 的刚性内型。从上面的讨论应当清楚，对我们最相关的是族 $\{G_b,\ b\in B(G)_{bsc}\}$，即 $G$ 的扩张纯内型。但还可以再走一步，一并考虑所有 $G_b$（$b\in B(G)$）的表示论。这些群的表示范畴\ysrc{106}粘合得很好，构成一个范畴；猜想上它同等于 $\Loc_{{}^{c}G}$ 上 (ind)-凝聚层的范畴，下面就来解释。

基本想法是：这些表示范畴粘成某个叠上的层范畴。事实上，单个 $\Rep(G_b(F),\Lambda)$ 可以看作（在适当意义下）局部紧群 $G_b(F)$ 的分类叠 $[*/G_b(F)]$ 上取 $\Lambda$-系数的层范畴。注意 $B(G)$ 是范畴 $\mathbf{Tor}_{G,\mathrm{iso}F}$（§4.2 引入）的底层，而每个 $b\in\mathbf{Tor}_{G,\mathrm{iso}F}$ 的自同构群是 $G_b(F)$。于是自然期望：$B(G)$ 是某个叠的 $\kappa_F$-点集，该叠在 $b$ 处的自同构群 $\mathrm{Aut}_b$ 是 $G_b(F)$（或与之密切相关的群），而所要的粘合范畴就是该叠上适当意义下的层范畴 $\mathrm{Shv}(\mathfrak B(G),\Lambda)$。特别地，对每个 $b\in B(G)$，应存在一对伴随函子
\begin{equation}
i_{b,!}:\Rep(G_b(F),\Lambda)\cong\mathrm{Shv}([*/{\mathrm{Aut}_b}],\Lambda)\rightleftarrows\mathrm{Shv}(\mathfrak B(G),\Lambda):i^!_b \tag{4.19}\label{eq:a2-4-19}
\end{equation}
其中 $i_b:[*/{\mathrm{Aut}_b}]\to\mathfrak B(G)$ 是相应的嵌入。

据我们所知，把这个想法变精确有两条路。一条属于 Fargues–Scholze。这条进路把 $B(G)$ 看作 Fargues–Fontaine 曲线上 $G$-丛的 $v$-叠 $\Bun_G$ 的点集，并把 $\mathrm{Shv}(B(G),\Lambda)$ 定义为 $\Bun_G$ 上适当定义的 étale 层范畴 $D(\Bun_G,\Lambda)$ [27]。这种定义相当精细，依赖 Scholze 关于菱形与凝聚数学的 $\ell$-进形式化的工作。

另一条进路\footnote{这一进路在几何 Langlands 圈内已流传了一段时间。}也许不那么精细，并且停留在概形的传统 $\ell$-进形式化\footnote{但对某些目的，例如 $p$-进局部 Langlands 纲领，这一进路可能不够。}的框架内：把 $B(G)$ 看作商叠
\[
\mathfrak B(G):=LG/\mathrm{Ad}_\sigma LG
\]
的点集，其中 $LG$ 记 $G$ 的圈群，即 $\kappa_F$ 上的（完美）群 ind-概形；$\mathrm{Ad}_\sigma$ 记 Frobenius 扭共轭，即 $\mathrm{Ad}_\sigma:LG\times LG\to LG$，$(h,g)\mapsto hg\sigma(h)^{-1}$（例如综述见 [86, 2.1]）。此时把 $\mathrm{Shv}(B(G),\Lambda)$ 定义为适当意义下的 $\Lambda$-层范畴 $\mathrm{Shv}(\mathfrak B(G)_{\bar\kappa_F},\Lambda)$。

更精确地说，这个范畴也可以（借助「h-下降」）实现为局部 shtuka 模叠 $\mathrm{Sht}^{\mathrm{loc}}$（腿在闭点 $0\in\Spec\mathcal O_F$ 处）上、态射由上同调对应给出的层范畴。[86] 末尾有简略讨论（另见 [32, 4.1]），这个范畴的详细研究将出现在 [88] 中。这里我们重复 [86] 中的梗概。以下所有几何对象都定义在 $\bar\kappa_F$ 上，尽管其中有些本来定义在 $\kappa_F$ 上。

先考虑一个较简单的情形，以定义 $G(F)$ 的分类叠上的 $\infty$-范畴 $\mathrm{Shv}([*/G(F)],\Lambda)$，它同等于 $G(F)$ 的光滑表示范畴 $\Rep(G(F),\Lambda)$。设 $K\subset G(F)$ 是开紧子群。由于可以把 $K$ 写成 $K=\varprojlim K_i$（其中每个 $K_i$ 有限），我们可以把 $\Lambda$ 看作 $\kappa_F$ 上的仿射群概形。考虑叠的群胚
\[
K\backslash G(F)/K\cong[*/K]\times_{[*/G(F)]}[*/K]\rightrightarrows[*/K],
\]
它延拓为单纯叠图（略去退化映射）
\begin{equation}
\cdots\rightrightarrows\rightrightarrows K\backslash G(F)/K\times_{[*/K]}K\backslash G(F)/K\rightrightarrows\rightrightarrows K\backslash G(F)/K\Rightarrow[*/K]. \tag{4.20}\label{eq:a2-4-20}
\end{equation}

\ysrc{107}
虽然 $[*/K]$ 与 $K\backslash G(F)/K$（以及上图的每一项）不是代数的，但它们仍可由 $\kappa_F$ 上有限型的好 (perfect) Deligne–Mumford 叠逼近，于是可以给它们附上 $\Lambda$-层的 $\infty$-范畴 $\mathrm{Shv}(-,\Lambda)$。例如可以定义 $\mathrm{Shv}([*/K],\Lambda)=\varinjlim\mathrm{Shv}([*/K_i],\Lambda)$，连接函子取沿有限群分类叠 $[*/K_i]\to[*/K_j]$ 的层拉回。则 $\mathrm{Shv}([*/K],\Lambda)=\Rep(K,\Lambda)$。对 $K\backslash G(F)/K$，可以把 $G(F)$ 写成 $K\times K$-稳定子集的递增并 $G(F)=\varinjlim_i G(F)_i$（从而把 $G(F)$ 看作 $\kappa_F$ 上的 ind-概形）。然后按上面的方式先定义范畴 $\mathrm{Shv}(K\backslash G(F)_i/K,\Lambda)$，再定义 $\mathrm{Shv}(K\backslash G(F)/K,\Lambda)=\varinjlim_i\mathrm{Shv}(K\backslash G(F)_i/K,\Lambda)$。

上图中所有态射都是 ind-可表示的（事实上 ind-有限）。于是可以把 $\mathrm{Shv}([*/G(F)],\Lambda)$ 定义为单纯 $\infty$-范畴
\[
\cdots\rightrightarrows\rightrightarrows\mathrm{Shv}(K\backslash G(F)/K\times_{[*/K]}K\backslash G(F)/K,\Lambda)\rightrightarrows\rightrightarrows\mathrm{Shv}(K\backslash G(F)/K,\Lambda)\Rightarrow\mathrm{Shv}([*/K],\Lambda)
\]
的几何实现，其中连接函子由真前推给出（[86, Remark 6.2]）。可以证明，这样定义的 $\mathrm{Shv}([*/G(F)],\Lambda)$ 不依赖 $\Lambda$ 的选取，并且确实同等于 $\Rep(G(F),\Lambda)$。

为定义 $\mathrm{Shv}(\mathfrak B(G),\Lambda)$，我们采用同样的策略，把 $\Lambda$ 换成 $G$ 在 $\mathcal O_F$ 上 Iwahori 模型 $\mathcal G$ 的正圈群 $L^+\mathcal G$（事实上也可以用 $G$ 的任意 parahoric 模型），并把 $[*/K]$ 换成
\begin{equation}
\mathrm{Sht}^{\mathrm{loc}}:=\frac{LG}{\mathrm{Ad}_\sigma L^+\mathcal G}, \tag{4.21}\label{eq:a2-4-21}
\end{equation}
即局部 $\mathcal G$-shtuka 的模叠（腿在 $0\in\Spec\mathcal O_F$，见 [86, (4.1.1)]）。再令
\begin{equation}
\mathrm{Hk}(\mathrm{Sht}^{\mathrm{loc}}):=\mathrm{Sht}^{\mathrm{loc}}\times_{\mathfrak B(G)}\mathrm{Sht}^{\mathrm{loc}} \tag{4.22}\label{eq:a2-4-22}
\end{equation}
为局部 shtuka 的 Hecke 叠（见 [86, (4.1.2)]，那里取 $s=t=1$）。类似地有单纯图
\begin{equation}
\cdots\rightrightarrows\rightrightarrows\mathrm{Hk}(\mathrm{Sht}^{\mathrm{loc}})\times_{\mathrm{Sht}^{\mathrm{loc}}}\mathrm{Hk}(\mathrm{Sht}^{\mathrm{loc}})\rightrightarrows\rightrightarrows\mathrm{Hk}(\mathrm{Sht}^{\mathrm{loc}})\Rightarrow\mathrm{Sht}^{\mathrm{loc}} \tag{4.23}\label{eq:a2-4-23}
\end{equation}
其中的态射是 ind-(perfectly) proper 的。同样，上图的每一项不是代数的，但可由 $\kappa_F$ 上有限型的好 (perfect) 代数叠逼近（详细讨论见 [77]，概述见 [86, 4.1]）。于是可以给每一项附上 $\Lambda$-层的 $\infty$-范畴，并把 $\mathrm{Shv}(\mathfrak B(G),\Lambda)$ 定义为相应单纯 $\infty$-范畴的几何实现。按定义，有自然函子 $\mathrm{Shv}(\mathrm{Sht}^{\mathrm{loc}},\Lambda)\to\mathrm{Shv}(\mathfrak B(G),\Lambda)$；这无非是沿 Newton 映射 $\mathrm{Nt}:\mathrm{Sht}^{\mathrm{loc}}\to\mathfrak B(G)$ 的真前推。

单纯图 (4.20) 经闭嵌入进入 (4.23)，该嵌入由
\begin{equation}
[*/I]\cong\frac{L^+\mathcal G}{\mathrm{Ad}_\sigma L^+\mathcal G}\subset\mathrm{Sht}^{\mathrm{loc}} \tag{4.24}\label{eq:a2-4-24}
\end{equation}
诱导，其中 $I=\mathcal G(\mathcal O_F)$。由此得到全忠实嵌入
\[
i_!:\Rep(G(F),\Lambda)\cong\mathrm{Shv}([*/G(F)],\Lambda)\to\mathrm{Shv}(\mathfrak B(G),\Lambda).
\]

\ysrc{108}
于是对每个开紧子群 $K'$，对象 $\delta_{K'}\in\Rep(G(F),\Lambda)$ 给出 $\mathrm{Shv}(\mathfrak B(G),\Lambda)$ 中的相应对象，仍记作同一记号。若 $K'\subset I$，则 $\delta_{K'}$ 在几何上由沿态射 $[*/K']\to[*/I]\to\mathrm{Sht}^{\mathrm{loc}}\to\mathfrak B(G)$ 对常值层 $\Lambda$ 作真前推给出。

\begin{rem}{4.34}\label{thm:a2-4-34}
如 [86] 所解释，$\mathrm{Shv}(\mathfrak B(G),\Lambda)$ 的同伦范畴可以表为 $\mathrm{Sht}^{\mathrm{loc}}$ 上的层范畴，其态射由支撑在 $\mathrm{Hk}(\mathrm{Sht}^{\mathrm{loc}})$ 上的上同调对应给出。后者已由 [77] 详细构造，对整体应用很有用。用这种解释，有更初等的办法看出：层 $\delta_{K'}$ 的自同态代数（定义为沿 $[*/K']\to\mathfrak B(G)$ 对 $\Lambda$ 作真前推）就是导出 Hecke 代数 $H_K$（见 [77, Remark 5.4.5]）。
\end{rem}

更一般地，对 basic 的 $b$，把它提升为元素 $\tilde b\in G(\breve F)$，后者正规化 $\mathcal G(\mathcal O_{\breve F})$；这里如以前一样，$\breve F$ 记 $F$ 的最大非分歧扩域的完备化。有与 (4.24) 类似的闭嵌入
\begin{equation}
[*/I_b]\cong\frac{L^+\mathcal G\cdot\tilde b}{\mathrm{Ad}_\sigma L^+\mathcal G}\subset\mathrm{Sht}^{\mathrm{loc}}. \tag{4.25}\label{eq:a2-4-25}
\end{equation}
这里 $I_b$ 是 $\tilde b$ 在 $\mathcal G(\mathcal O_{\breve F})$ 中的扭中心化子，它是 $G_b(F)$ 的 Iwahori 子群。于是有与 (4.20) 类似、附于群胚 $[*/I_b]\times_{[*/G_b(F)]}[*/I_b]\rightrightarrows[*/I_b]$ 的单纯图，并且它闭嵌入 (4.23)。这就给出 (4.19) 中所承诺的嵌入 $i_{b,!}$。

\begin{rem}{4.35}\label{thm:a2-4-35}
最佳的猜测是：Fargues–Scholze 定义的范畴 $D(\Bun_G,\Lambda)$ 与上面概述的 $\mathrm{Shv}(\mathfrak B(G),\Lambda)$ 等价。一个显著的现象是：在 $B(G)$ 的上述两种解释中，$B(G)$ 上的偏序被反转。
\end{rem}

\begin{rem}{4.36}\label{thm:a2-4-36}
如 [86] 所提到，完全同样的构造可以用来定义并研究旁伴随商空间 $LG/\mathrm{Ad}\,LG$ 上的层范畴。
\end{rem}

现在来表述我们的猜想。令 $\hat{\mathcal N}_{{}^{c}G}$ 记 $\Sing(\Loc_{{}^{c}G})$ 的子集，如 (3.13) 所述。$\Loc_{{}^{c}G}$ 上凝聚层范畴的约定见注 3.8。

\begin{conj}{4.37}\label{thm:a2-4-37}
设 $(G,B,T,e)$ 是 $F$ 上带 pinning 的拟分裂约化群。则存在自然的 $\mathbf{Tor}_{Z_G,\mathrm{iso}F}$-等变 $\infty$-范畴等价
\[
\mathbb L_G:\mathrm{Shv}(\mathfrak B(G),\Lambda)\to\mathrm{Ind}\,\Coh_{\hat{\mathcal N}_{{}^{c}G}}(\Loc_{{}^{c}G}),
\]
它把 $\mathrm{Whit}_{(U,\psi)}$（见 (4.5)）送到结构层 $\mathcal O_{\Loc_{{}^{c}G}}$。此外，对每个 basic 元 $b\in B(G)$，猜想 4.24 中的猜想函子 $\mathfrak A_{G_b}$ 张量 $\Lambda$ 后落入下述交换图：
\[
\begin{tikzcd}[column sep=large,row sep=large]
\Rep_{\mathrm{f.g.}}(G_b,\Lambda) \arrow[r, "\mathfrak A_{G_b}"] \arrow[d, "i_{b,!}"'] & \Coh(\Loc_{{}^{c}G}) \arrow[d] \\
\mathrm{Shv}(\mathfrak B(G),\Lambda) \arrow[r, "\mathbb L_G"'] & \mathrm{Ind}(\Coh_{\hat{\mathcal N}_{{}^{c}G}}(\Loc_{{}^{c}G}))
\end{tikzcd}
\]
\end{conj}

\ysrc{109}
\begin{rem}{4.38}\label{thm:a2-4-38}
注意该猜想蕴含：对每个 $b$（不必 basic），$\Loc_{{}^{c}G}$ 上应存在 ind-凝聚层
\[
\mathfrak A_{G_b,\{1\}}:=\mathbb L_G(i_{b,!}(\delta_{G_b,\{1\}})),\qquad \delta_{G_b,\{1\}}:=C_c(G_b(F),\Lambda),
\]
其中 $i_{b,!}$ 是 (4.19) 中的函子，$\nu_b$ 是 $b$ 的 Newton 点（即 $\hat G$ 的支配有理特征 [47, 4.2]）。一般说来，$\mathfrak A_{G_b,\{1\}}$ 不应是通常的凝聚层，这与猜想 4.24(2) 所考虑的 basic 情形不同。但我们猜想 $\mathfrak A_{G_b,\{1\}}$ 集中在上同调次数 $[-(2\rho,\nu_b),0]$。
\end{rem}

\begin{rem}{4.39}\label{thm:a2-4-39}
在 Fargues–Scholze 的进路中，$\mathrm{Shv}(B(G),\Z_\ell)$ 定义为 $D(\Bun_G,\Z_\ell)$，上述猜想在形式上看起来就像 Arinkin–Gaitsgory 提出的整体几何 Langlands 猜想 [2]。事实上，Fargues–Scholze 独立宣布了用 $D(\Bun_G,\Z_\ell)$ 表述的同一猜想。
\end{rem}

\begin{rem}{4.40}\label{thm:a2-4-40}
当系数取 $\Z_\ell$、且 $\ell$ 是所谓的非 banal 素数（non banal prime）时，$\mathfrak A_{G_b}$ 的存在性不能直接由 $\mathbb L_G$ 的存在性得出，因为 $\Rep_{\mathrm{f.g.}}(G_b,\Z_\ell)$ 并不落在 $\mathrm{Shv}(\mathfrak B(G),\Z_\ell)$ 的紧对象子范畴中。不过，有重整化的版本 $\mathrm{Ind}\,\mathrm{Shv}_{\mathrm{f.g.}}(\mathfrak B(G),\Z_\ell)$，其紧对象子范畴将包含 $\Rep_{\mathrm{f.g.}}(G_b,\Z_\ell)$（定义类似于 [2, 12.2.3]，将在 [88] 中给出）。我们预期 $\mathbb L_G$ 延拓为等价
\[
\mathbb L_G^{\mathrm{Ind\,f.g.}}:\mathrm{Ind}\,\mathrm{Shv}_{\mathrm{f.g.}}(\mathfrak B(G),\Z_\ell)\cong\mathrm{Ind}\,\Coh(\Loc_{{}^{c}G}),
\]
这将蕴含 $\mathfrak A_{G_b}$ 的存在性。注意把 $\Z_\ell$ 换成 $\Q_\ell$ 后，$\mathrm{Ind}\,\mathrm{Shv}_{\mathrm{f.g.}}(\mathfrak B(G),\Q_\ell)=\mathrm{Shv}(\mathfrak B(G),\Q_\ell)$，而幂零奇异支撑条件由引理 3.23 自动满足。因此 $\mathbb L_G^{\mathrm{Ind\,f.g.}}$ 与 $\mathbb L_G$ 一致。
\end{rem}

\begin{rem}{4.41}\label{thm:a2-4-41}
值得把上述等价的「motivic」（即不依赖 $\ell$ 的）版本表述出来。当系数 $\Lambda=\Q_\ell$ 时，命题 3.12 提示：在伽罗瓦侧可以不必考虑 $\Coh(\Loc_{{}^{c}G,F}\otimes\Q_\ell)$，而考虑 $\Coh(\Loc^{WD}_{{}^{c}G,F}/\mathbb G_m\otimes\Q_\ell)$。另一方面，我们预期 $\mathrm{Shv}(\mathfrak B(G),\Q_\ell)$ 容许混合版本 $\mathrm{Shv}^m(\mathfrak B(G),\Q_\ell)$。于是 $\mathbb L_G$ 也许能提升为混合范畴之间的等价，从而有机会下降到 $\Q$ 上。
\end{rem}

\begin{rem}{4.42}\label{thm:a2-4-42}
上述猜想等价应当满足一组与整体几何 Langlands 对应相类似的相容性条件（[2, 31]）。例如，它应同时与两侧的抛物诱导相容，也应与 $\mathrm{Shv}(\mathfrak B(G),\Lambda)$ 上的上同调对偶（(4.3) 的推广）以及修正的 Grothendieck–Serre 对偶 (3.9) 相容。由于讨论这些相容性需要引入与 $\mathrm{Shv}(\mathfrak B(G),\Lambda)$ 相关的更多构造，这里从略，详见 [88]。
\end{rem}

另一方面，上述猜想等价预测：$\Loc_{{}^{c}G}$ 上完美复形的范畴 $\Perf(\Loc_{{}^{c}G})$ 应作用在 $\mathrm{Shv}(\mathfrak B(G),\Lambda)$ 上，通常称之为谱作用（spectral action）。Fargues–Scholze 已宣布在他们设定下构造了这样的作用。但在 $\mathrm{Shv}(\mathfrak B(G),\Lambda)$ 上，这样的谱作用是否存在尚不清楚。

作为 $\mathrm{Shv}(\mathfrak B(G),\Z_\ell)$ 可能也是该猜想的正确输入的一个证据，我们先回忆 [77, 83, 86] 中的下述结果。

\ysrc{110}
\begin{thm}{4.43}\label{thm:a2-4-43}
设 $G$ 约化。则存在自然函子 $\Coh(\Loc^{\mathrm{ur}}_{{}^{c}G})\to\mathrm{Shv}(\mathfrak B(G),\Lambda)$，使下图交换：
\[
\begin{tikzcd}[column sep=large,row sep=large]
\Rep(\hat G,\Lambda)^{\heartsuit} \arrow[r, "\mathrm{Sat}"] \arrow[d] & \mathrm{Shv}(\mathrm{Sht}^{\mathrm{loc}},\Lambda) \arrow[d] \\
\Coh(\Loc^{\mathrm{ur}}_{{}^{c}G}) \arrow[r] & \mathrm{Shv}(\mathfrak B(G),\Lambda)
\end{tikzcd}
\]
其中 $\mathrm{Sat}$ 由几何 Satake 等价诱导（[60, 82, 85]），而左侧的竖直函子是沿 $\Loc^{\mathrm{ur}}_{{}^{c}G}\to\mathbb B\hat G$ 的自然拉回函子。
\end{thm}

更有说服力的是下述陈述，它将由 Hemo 与作者在 [88] 中建立。

\begin{thm}{4.44}\label{thm:a2-4-44}
设 $(G,B,T,e)$ 是局部域 $F$ 上带 pinning 的非分歧群，且 $\Lambda=\overline\Q_\ell$。则定理 4.43 中的函子延拓为全忠实嵌入
\[
\Coh(\Loc^{\widehat{\mathrm{unip}}}_{{}^{c}G})\into\mathrm{Shv}(\mathfrak B(G),\overline\Q_\ell),
\]
其像落在 $\mathrm{Shv}(\mathfrak B(G),\overline\Q_\ell)$ 的紧对象子范畴中。它把 $\Coh\mathrm{Spr}^{\mathrm{unip}}_{{}^{c}G}$ 送到 $\delta_I$。更一般地，对每个元 $b\in B(G)$，令 $H_{I_b}$ 是 $G_b$ 的相应 Iwahori–Hecke 代数。则有下述交换图：
\[
\begin{tikzcd}[column sep=large,row sep=large]
\mathrm{Mod}_{H_{I_b}} \arrow[d, hook] & \Rep(G_b(F),\overline\Q_\ell) \arrow[l] \arrow[d, "i_{b,!}"'] \\
\mathrm{Ind}\,\Coh(\Loc^{\widehat{\mathrm{unip}}}_{{}^{c}G}) \arrow[r] & \mathrm{Shv}(\mathfrak B(G),\overline\Q_\ell)
\end{tikzcd}
\]
\end{thm}

定理中嵌入的更多性质将在 [88] 中研究。

\subsection{模簇的上同调与局部–整体相容性}\label{sec:a2-4-7}

在本小节中，我们对 shtuka 模叠的上同调表述猜想性公式，并给出一些证据。我们主要考虑函数域情形，因为那里的图景更完整。但我们也会讨论通过志村簇的上同调实现 Jacquet–Langlands 传输的猜想性几何化，这是 [77] 主构造的推广。

设 $F$ 是整体域，$G$ 是 $F$ 上连通约化群。设 $\Lambda$ 是诺特 $\Z_\ell$-代数；若 $F$ 是函数域，则要求 $\ell\ne\operatorname{char}F$。我们沿用 §3.4 的记号。

先讨论函数域情形。设 $F=\F_q(X)$ 是整体函数域，其中 $X$ 是 $\F_q$ 上几何连通的光滑射影曲线。以 $W_F$ 记 $F$ 的 Weil 群。设 $\eta=\Spec F$ 是 $X$ 的一般点，$\bar\eta$ 是 $\eta$ 上的几何点。令 $\mathbb O=\prod_{v\in|X|}\mathcal O_v$ 为整数阿代尔，其中 $\mathcal O_v\subset F_v$ 是整数环。把 $G$ 延拓为 $X$ 上的 Bruhat–Tits 整模型 $\mathcal G$，即 $X$ 上的光滑仿射群概形，使得对每个 $v$，$\mathcal G|_{\mathcal O_v}$ 是 Bruhat–Tits 意义下 $G_v$ 的 parahoric 群概形。我们要考虑 $\mathcal G$-shtuka 模叠的紧支集上同调。$\mathcal G$-shtuka 模叠的基本构造与事实可参见 [49]。

固定一个 level $K\subset\mathcal G(\mathbb O)$。令 $S_K$ 是满足 $K_v\ne\mathcal G(\mathcal O_v)$ 的位点 $v$ 的集合，$S\supset S_K$ 是 $K_v$ 非超特殊的位点集合。对有限集 $I$，以 $\mathrm{Sht}_{(X-S_K)^I,K}$ 记 $X$ 上在 $X-S_K$ 中带 $I$ 条腿、\ysrc{111}带 $K$-level 结构的 $\mathcal G$-shtuka 模叠。它是 $(X-S_K)^I$ 上的 ind-Deligne–Mumford 叠。沿对角映射 $\bar\eta\to(X-S_K)^I$ 对它作底变，所得叠记作 $\mathrm{Sht}_{\Delta(\bar\eta),K}$。对 $({}^{c}G)^I$ 在有限投射 $\Lambda$-模上的每个表示 $V$，几何 Satake 对应给出 $\mathrm{Sht}_{\Delta(\bar\eta),K}$ 上的反常层 $\mathrm{Sat}(V)$（其实它定义在 $\mathrm{Sht}_{(X-S)^I,K}$ 上）。令
\[
C_c(\mathrm{Sht}_{\Delta(\bar\eta),K},\mathrm{Sat}(V))\in\mathbf{Mod}_{H_K}
\]
$C_c(\mathrm{Sht}_{\Delta(\bar\eta),K},\mathrm{Sat}(V))$ 记以 $\mathrm{Sat}(V)$ 为系数的总紧支集上同调（即其上链复形），相应的整体（导出）Hecke 代数（系数在 $\Lambda$ 中）$H_K=C_c(K\backslash G(\mathbb A)/K,\Lambda)$ 作用在其上。当 $V=\mathbf 1$ 是平凡表示时，有
\[
C_c(\mathrm{Sht}_{\Delta(\bar\eta),K},\mathrm{Sat}(\mathbf 1))=\coprod_{\xi\in\ker^1(F,G)}C_c(G^\xi(F)\backslash G(\mathbb A)/K,\Lambda).
\]
这里 $\ker^1(F,G)\subset H^1(F,G)$ 由局部平凡的那些类组成；对 $\xi\in\ker^1(F,G)$，$G^\xi$ 记相应的纯内型；$G^\xi(F)\backslash G(\mathbb A)/K$ 看作 $\bar\eta$ 上的离散 DM 叠，而 $C_c(G^\xi(F)\backslash G(\mathbb A)/K,\Lambda)$ 记其紧支集上同调。当 $\Lambda=\Q_\ell$ 且 $G$ 满足 Hasse 原理（例如 $G$ 拟分裂）时，这就是 $G(F)\backslash G(\mathbb A)/K$ 上紧支集函数构成的空间。

令 $H^i_{I,V}=H^iC_c(\mathrm{Sht}_{\Delta(\bar\eta),K},\mathrm{Sat}(V))$。由 [79–81]，自然的伽罗瓦作用与部分 Frobenius 作用合起来诱导出 $H^i_{I,V}$ 上的典范 $W^I_{F,S}$-作用。下述陈述可以看作 [50] 主构造的推广。

\begin{thm}{4.45}\label{thm:a2-4-45}
设 $\Lambda=\Q_\ell$，并把 $\Loc_{{}^{c}G,F,S}$ 看作 $\Q_\ell$ 上的代数叠。则对每个 $i$，存在 ${}^{cl}\Loc_{{}^{c}G,F,S}$ 上的拟凝聚层 $A^i_\Lambda$，带有 $H_K$ 的作用，使得对 $({}^{c}G)^I$ 的每个有限维表示 $V$，有自然的 $(H_K\times W^I_{F,S})$-等变同构
\begin{equation}
H^i_{I,V}\cong\Gamma({}^{cl}\Loc_{{}^{c}G,F,S},(\mathcal W_{F,S}V)\otimes A^i_\Lambda), \tag{4.26}\label{eq:a2-4-26}
\end{equation}
其中 $\mathcal W_{F,S}V$ 是 $\Loc_{{}^{c}G,F,S}$ 上带 $W^I_{F,S}$ 作用的向量丛，见注 2.13。
\end{thm}

\begin{proof}
如 [50, §5] 所解释，对 $\hat G\times({}^{c}G)^I$ 的表示 $V$，可以定义 $H^i_{\{0\}\cup I,V}$，它带有 $H_K\times W^I_F$ 的作用，并且当 $V$ 在 $\hat G$ 因子上的限制平凡时 $H^i_{\{0\}\cup I,V}=H^i_{I,V}$。特别地，有 $H_K$-模 $H^i_{\{0\},\mathrm{Reg}}$，其中 $\mathrm{Reg}$ 记 $\hat G$ 的正则表示。

把 $W_{F,S}$ 看作抽象群，并考虑 ${}^{cl}\mathcal R_{W_{F,S},{}^{c}G}$。[50, §6] 的构造给出同态 $\Q_\ell[{}^{cl}\mathcal R_{W_{F,S},{}^{c}G}]\to\End(H^i_{\{0\},\mathrm{Reg}})$。令 $A^i$ 是该映射的像。对 $f\in\Q_\ell[{}^{c}G]$ 与 $\gamma\in W_{F,S}$，有 ${}^{cl}\mathcal R_{W_{F,S},{}^{c}G}$ 上的正则函数 $F_{f,\gamma}$，它由 $F_{f,\gamma}(\rho)=f(\rho(\gamma))$ 给出。令 $\bar F^i_{f,\gamma}$ 是 $F_{f,\gamma}$ 在 $A^i$ 中的像。注意，若把 $H^i_{I,V}$ 看作 $\pi_1(\overline Y)^I$ 的表示，则它是有限维连续子表示的并。于是 [50, 6.2] 与引理 2.57 中的论证表明：映射 $\pi_1(\overline Y)\to A^i$，$\gamma\mapsto\bar F_{f,\gamma}$ 连续，这里 $A^i$ 配备 ind-$\ell$-adic 拓扑。因此有分解
\[
\Spec A^i\to{}^{cl}\mathcal R^{sc}_{W_{F,S},{}^{c}G}\to{}^{cl}\mathcal R_{W_{F,S},{}^{c}G}.
\]
\end{proof}

\ysrc{112}
于是 $H^i_{\{0\},\mathrm{Reg}}$ 可以看作 ${}^{cl}\mathcal R^{sc}_{W_{F,S},{}^{c}G}$ 上的拟凝聚层。如 [50] 所解释，$H^i_{\{0\},\mathrm{Reg}}$ 上还有与 $A^i$ 的作用相容的 $\hat G$-作用，故它下降为 ${}^{cl}\mathcal R^{sc}_{W_{F,S},{}^{c}G}/\hat G$ 上的拟凝聚层 $\mathfrak A^i_\Lambda$。由构造，$\mathfrak A^i_\Lambda$ 支撑在 ${}^{cl}\Loc_{{}^{c}G,F,S}$ 上；再由 [50] 中的论证即知 (4.26) 成立。

\begin{rem}{4.46}\label{thm:a2-4-46}
如 [50] 所解释，层 $\mathfrak A^i_\Lambda$ 其实是 $({}^{cl}\Loc^{\square}_{{}^{c}G,F,S}/(\hat G/Z^{\Gamma_F}_{\hat G}))\otimes\Q_\ell$ 上某个拟凝聚层的拉回。我们预期每个 $\mathfrak A^i_\Lambda$ 都是凝聚的。
\end{rem}

\begin{exam}{4.47}\label{thm:a2-4-47}
设 $G$ 半单（为简单起见），并回忆例 3.33 中的椭圆 Langlands 参数。由此可知：$\mathfrak A^i_\Lambda$ 在椭圆参数 $\rho$ 处的局部化（记作 $\mathfrak A^i_{K,\rho}$）是带 $H_K\times S_\rho$ 作用的 $\overline\Q_\ell$-向量空间。而 $H^i_{I,W}$ 在 $\rho$ 处的局部化同构于 $(\mathfrak A^i_{K,\rho}\otimes W_\rho)^{S_\rho}$。\tn{此处 $H^i_{I,W}$ 的下标与定理 4.45 中的 $H^i_{I,V}$ 不一致，且 $W$ 未加定义，疑为 $H^i_{I,V}$ 之误；译文照原样保留。}因此定理 4.45 恢复了 [50] 的主要结果（除 $\mathfrak A^i_{K,\rho}$ 的有限维性之外）。该公式与 Arthur–Kottwitz 重数公式之间的关系，见 loc. cit.。
\end{exam}

\begin{rem}{4.48}\label{thm:a2-4-48}
(1) 应有形如 (4.26) 的陈述，这一想法源自 Drinfeld，是对 [49] 中某个构造的解释。如 [3, 32, 33] 所解释，(4.26) 的（导出版本）应当由范畴几何 Langlands 对应取范畴迹得到。

(2) 当 $\Lambda=\Z_\ell$ 时，一般不应期待定理 4.45 成立。问题在于：函子 $V\mapsto H^i_{I,V}$ 与函子 $\Gamma({}^{cl}\Loc_{{}^{c}G,F,S},-)$ 对整系数都不 $t$-正合。不过我们确实期待：把公式中的各个上同调群合成总上链复形 $C_c(\mathrm{Sht}_{\Delta(\bar\eta),K},\mathrm{Sat}(V))$、把各个 $\mathfrak A^i_\Lambda$ 合成 $\Loc_{{}^{c}G,F,S}$ 上的拟凝聚复形后，(4.26) 的导出版本成立。精确的猜想见下。
\end{rem}

在 [50] 中，受 Arthur–Kottwitz 猜想启发，我们猜想 $\mathfrak A^i_\Lambda$ 分解为局部因子的张量积。现在我们进一步猜想：这些局部因子恰好就是猜想 4.24 中出现的那些凝聚层。为简单起见，本小节余下部分假定 $G$ 的中心 $Z_G$ 连通。

为精确表述猜想，先注意在整体情形下也可以类似地定义 $\mathbf W_G$、$\mathbf{Tor}_{G,\mathrm{iso}F}$ 与 $\mathbf{TS}_G$（§4.2 引入）：只须在相应的构造中，把那里局部域的最大非分歧扩域的完备化换成整体情形中 $F$ 的最大非分歧扩域。$\mathbf{Tor}_{G,\mathrm{iso}F}$ 的同构类集合仍记作 $B(G)$，basic 元构成的子集 $B(G)_{bsc}$ 也类似地定义。$G$ 的一个整体 basic 元在每个位点给出 $G_v$ 的局部 basic 元（其在 $\mathbb X^{\bullet}(Z^{\Gamma_v}_{\hat G})$ 中的像对几乎所有 $v$ 为零），并有下述带点集的正合列：
\[
B(G)_{bsc}\to\bigoplus_v B(G_v)_{bsc}\to\mathbb X^{\bullet}(Z^{\Gamma_F}_{\hat G}).
\]

现在固定一个非平凡特征 $\psi_0:F\backslash\mathbb A\to\Lambda^\times$，并固定一个整体元 $t\in\mathbf{TS}_G$。这些数据在每个局部位点诱导出相应的数据。于是\ysrc{113}每个位点 $v$ 处都有猜想 4.24 中的函子 $\mathfrak A_{G_v}$。若 $K_v\subset G_v$ 是开紧子群，为简单起见我们有时把 $\mathfrak A_{G_v,K_v}$ 简写作 $\mathfrak A_{K_v}$。

回忆我们固定了 level 结构 $\Lambda$。必要时放大集合 $S$，可以假设对每个 $v\notin S$ 都有 $t_v\in\mathbf W_{G_v}$，且 $K_v$ 是由 pinning 决定的超特殊子群（至 $G(F_v)$-共轭）。以 $\boxtimes_{v\in S}\mathfrak A_{K_v}$ 记 $\prod_{v\in S}\Loc_v$ 上这些凝聚层的外张量积，以 $\mathrm{res}^!(\boxtimes_{v\in S}\mathfrak A_{K_v})$ 记它沿 (3.27) 到 $\Loc_{{}^{c}G,F,S}$ 的 !-拉回。按我们的预期 (4.16)，$\mathrm{res}^!(\boxtimes_{v\in S}\mathfrak A_{K_v})$ 应不依赖 $t\in\mathbf{TS}_G$（以及 $\psi_0$）的选取，并且下降到 $\Loc^{\square}_{{}^{c}G,F,S}/(\hat G/Z^{\Gamma_F}_{\hat G})$ 上的拟凝聚层。

\begin{conj}{4.49}\label{thm:a2-4-49}
对 $({}^{c}G)^I$ 在自由 $\Lambda$-模上的每个表示 $V$，存在典范的 $(H_K\times W^I_{F,S})$-等变同构
\[
C_c(\mathrm{Sht}_{\Delta(\bar\eta),K},\mathrm{Sat}(V))\cong\Gamma(\Loc_{{}^{c}G,F,S},(\mathcal W_{F,S}V)\otimes\mathrm{res}^!(\boxtimes_{v\in S}\mathfrak A_{K_v})).
\]
\end{conj}

注意该猜想与放大 $S$ 是相容的：当 $K_w$ 超特殊（且由 $t_w$ 决定）时 $\mathfrak A_{K_w}\cong\mathcal O_{\Loc^{\mathrm{unr}}_w}$，而由引理 3.34，相应的图是笛卡儿方图。

\begin{rem}{4.50}\label{thm:a2-4-50}
为简单起见，设 $G$ 是旁伴随型。设 $\rho$ 是例 3.33 那样的椭圆 Langlands 参数。由于 $\rho$ 是孤立光滑点，$(\mathcal W_{F,S}V)\otimes\mathrm{res}^!(\boxtimes_{v\in S}\mathfrak A_{K_v})$ 在 $\rho$ 处的局部化是由 $V\otimes(\otimes_{v\in S}A^!_{K,v})$ 给出的向量空间复形，其中 $A^!_{K,v}$ 记 $\mathfrak A_{K,v}$ 在 $\rho_v:=\rho|_{W_v}$ 处的 !-纤维。由于 $\mathrm{Ad}_\rho$ 是权为 $0$ 的纯的，每个 $\rho_v$ 都是 $\Loc_v$ 的光滑点（命题 3.24）。而 $\mathfrak A_{K_v}$ 应是极大的 Cohen–Macaulay 通常凝聚层（猜想 4.24(2)）。这将蕴含 $A^!_{K,v}$ 位于上同调次数零。由此可见：例 4.47 中的 $\mathfrak A^i_{K,\rho}$ 当 $i\ne0$ 时应消失。这与一般的预期一致。
\end{rem}

\begin{exam}{4.51}\label{thm:a2-4-51}
我们把该猜想在处处非分歧的情形写得更具体，即 $\mathcal G$ 在 $X$ 上约化、$K=\mathcal G(\mathbb O)$。此时可考虑 $\Loc_{{}^{c}G,X}=\Loc_{{}^{c}G,F,\emptyset}$，如注 3.36 所述。由于 $\mathfrak A_{K_v}\cong\mathcal O_{\Loc^{\mathrm{unr}}_v}\cong\omega_{\Loc^{\mathrm{unr}}_v}$，猜想 4.49 在此情形化为
\[
C_c(\mathrm{Sht}_{\Delta(\bar\eta),K},\mathrm{Sat}(V))\cong\Gamma(\Loc_{{}^{c}G,X},(\mathcal W_FV)\otimes\omega_{\Loc_{{}^{c}G,X}}).
\]
我们指出，当 $G$ 分裂且 $\Lambda=\overline\Q_\ell$ 时，[3] 也独立地猜想了这一公式。\footnote{只是 loc. cit. 中 $\Loc_{{}^{c}G,X}$ 的定义先验地不同。}
\end{exam}

我们进一步特殊化到 $X=\mathbb P^1$、$V=\mathbf 1$（${}^{c}G$ 的平凡表示）的情形。此时 $G$ 必拟分裂，并且在常数域的扩张 $\F_{q'}/\F_q$ 上分裂。于是如前面提到的，$C_c(\mathrm{Sht}_{\Delta(\bar\eta),K},\mathrm{Sat}(\mathbf 1))$ 就是 $G(F)\backslash G(\mathbb A)/G(\mathbb O)$（看作离散 DM 叠）的紧支集上同调。若 $\Lambda=\Q_\ell$，这就是 $G(F)\backslash G(\mathbb A)/G(\mathbb O)$ 上紧支集函数构成的空间。

把双陪集 $G(F)\backslash G(F)G(\mathbb O)/G(\mathbb O)$ 的特征函数看作映射 $k\to C_c(G(F)\backslash G(\mathbb A)/G(\mathbb O),\Lambda)$。导出 Hecke 代数 $H_{K_0}=\End\delta_{K_0}$（$0\in\mathbb P^1$ 处）在 $H_\emptyset(\mathbf 1)=C_c(G(F)\backslash G(\mathbb A)/G(\mathbb O),\Lambda)$ 上的作用，诱导出 Radon 变换的一个导出版本
\[
H_{K_0}\cong H_{K_0}\otimes k\to H_{K_0}\otimes C_c(G(F)\backslash G(\mathbb A)/G(\mathbb O),\Lambda)\to C_c(G(F)\backslash G(\mathbb A)/G(\mathbb O),\Lambda),
\]
